\documentclass{template}

\usepackage{mathtools}
\pgfplotsset{compat=1.18}

\newcommand{\qedwhite}{\hfill \ensuremath{\Box}}


\begin{document}

\begin{center}
    \Large \textcolor{blue!70!black}{\textbf{Dynamical Systems Homework \#1}} \\
    \normalsize \textit{Dmytro Zakharov}
\end{center}

\begin{problem}[Exercise 1]
    Consider the following dynamical systems:
    \begin{enumerate}
        \item[(i)] $X=\mathbb{R}^2/\mathbb{Z}^2$, $T(x,y) = (x+\alpha,y+\beta)$ (for $\alpha,\beta \in \mathbb{R}/\mathbb{Z}$);
        \item[(ii)] $X=\mathbb{R}^2/\mathbb{Z}^2$, $T(x,y) = (ax, by)$ (for $a,b \in \mathbb{Z}$). 
    \end{enumerate}
    For each system, determine (depending on the parameters):
    \begin{enumerate}
        \item[(a)] Are there any fixed points?
        \item[(b)] Are there any points with least period 2?
        \item[(c)] Is the system minimal?
        \item[(d)] Is the system transitive?
        \item[(e)] Is the system topologically ergodic?
        \item[(f)] Is the system topologically weakly mixing?
        \item[(g)] Is the system topologically strongly mixing?
        \item[(h)] Is the system equicontinuous?       
    \end{enumerate}
\end{problem}

\textbf{Solution.} 

\textcolor{blue!70!black}{\textbf{Part \textit{(i)}.}}

\textit{Part (a).} Fixed point $(x_0,y_0) \in \mathbb{R}^2/\mathbb{Z}^2$ must
satisfy $x_0+\alpha=x_0$ and $y_0+\beta=y_0$. If $\alpha,\beta \in \mathbb{Z}$,
then any point is a fixed point (in fact, the map $T$ is such case is simply an
identity map). Otherwise, $T$ has no fixed points.

\underline{Further assume $\alpha,\beta \not\in \mathbb{Z}$, otherwise the map
$T$ is trivial.}

\textit{Part (b).} Note than $T^n(x,y) = (x+n\alpha,y+n\beta)$. If the point
$(x_0,y_0)$ is a point with least period $2$, then by definition $T^2(x_0,y_0) =
(x_0,y_0)$ (note that the point cannot have a period $1$ since we assumed
$\alpha,\beta \not\in \mathbb{Z}$), which translates to $2\alpha=0$ and
$2\beta=0$. This means that \textcolor{green!50!black}{$(\alpha,\beta) \in
\{0,1/2\}$}. Otherwise, $T$ has no points with least period $2$.

\textit{Part (c).} The system is minimal iff for any $(x,y) \in
\mathbb{R}^2/\mathbb{Z}^2$ we have $X=\overline{\mathcal{O}(x,y)}$. Note that is
suffices to consider when the orbit of point $(0,0)$ is dense (since for any
other point $(a,b) \in \mathbb{R}^2/\mathbb{Z}^2$, the homeomorphism
$\varphi_{a,b}: (x,y) \mapsto (x+a,y+b)$ maps the orbit of $(0,0)$ onto the
orbit of $(a,b)$, thus maps dense orbit to a dense orbit). The orbit of $(0,0)$
is precisely $\{(n\alpha,n\beta): n \in \mathbb{Z}\}$. We claim that this set is
dense iff \textcolor{green!50!black}{$\{1,\alpha,\beta\}$ are linearly
independent over $\mathbb{Q}$} (e.g., there is no $n,m \in \mathbb{Z}\setminus
\{0\}$ such that $n\alpha+m\beta \in \mathbb{Z}$). Indeed, irrationality of both
$\alpha$ and $\beta$ is required since otherwise one of the components will not
be dense. However, this condition is not sufficient (take, e.g., $\alpha :=
\sqrt{2}$ and $\beta := 2\sqrt{2}$): we need to ensure that we do not get the
closed curve. This can only happen if $\alpha/\beta \in \mathbb{Q}$.

\textit{Part (d).} If $\{1,\alpha,\beta\}$ are linearly independent over
$\mathbb{Q}$, then as we have just shown, $(X,T)$ is minimal and thus is
transitive and topologically ergodic. If these elements were not linearly
independent, then as we have shown, there would be no dense orbits, and thus the
system cannot be transitive. 

\textit{Part (e--g).} Since topological ergodicity is equivalent to
transitivity, the answer is the same as in \textit{(d)}. 

As we will see, the system is equicontinuous and thus is not topologically
weakly mixing (and consequently not strongly mixing) for any choice of
$\alpha,\beta$.

\textit{Part (h).} Recall that the definition of $(X,T)$ being equicontinuous
is as follows:
\begin{equation*}
    \forall_{(x_0,y_0) \in \mathbb{R}^2/\mathbb{Z}^2} \; \forall_{\varepsilon>0} 
    \; \exists_{U \ni (x_0,y_0)} \; d(T^n(x_0,y_0), T^n(x_1,y_1)) < \varepsilon \; \forall_{(x_1,y_1) \in U} \; \forall_{n \in \mathbb{Z}_{\geq 0}}
\end{equation*} 

Since our space is a compact metric space,
\begin{equation*}
    \forall_{\varepsilon > 0} \; \exists_{\delta > 0} \; d((x_0,y_0),(x_1,y_1))<\delta \Rightarrow d(T^n(x_0,y_0), T^n(x_1,y_1))<\varepsilon \; \forall_{n \in \mathbb{Z}_{\geq 0}}
\end{equation*}

But note that since translation is a distance-preserving function,
$d((x_0,y_0),(x_1,y_1))=d(T^n(x_0,y_0),T^n(x_1,y_1))$, thus this condition is trivially
satisfied (for $\delta := \varepsilon$). Thus $(X,T)$ is equicontinuous
\textcolor{green!50!black}{for any values of $\alpha,\beta$.}

\textcolor{blue!70!black}{\textbf{Part \textit{(ii)}.}}

\textit{Part (a).} Fixed point $(x_0,y_0) \in \mathbb{R}^2/\mathbb{Z}^2$ must
satisfy $ax_0=x_0$ and $by_0=y_0$, which is equivalent to requiring $(a-1)x_0
\in \mathbb{Z}$ and $(b-1)y_0 \in \mathbb{Z}$. Note that for either $a=1$ or
$b=1$, $T$ acts as an identity to the corresponding component. Therefore, we 
have four cases:
\begin{enumerate}
    \item[(a)] if $a=b=1$, $T$ is an identity, so all points are fixed;
    \item[(b)] if $a=1,b\neq 1$, then fixed points are $(\mathbb{R}/\mathbb{Z})
    \times \{\frac{m}{b-1}: m \in \mathbb{Z}\}$; there are in total
    $|b-1|$ such fixed points;
    \item[(c)] if $a \neq 1, b = 1$, then fixed points are $\{\frac{n}{a-1}: n \in \mathbb{Z}\} \times (\mathbb{R}/\mathbb{Z})$;
    there are in total $|a-1|$ such fixed points;
    \item[(d)] if $a,b \neq 1$, then fixed points are $\{\frac{n}{a-1},\frac{m}{b-1}: n,m
    \in \mathbb{Z}\}$; there are in total $|(a-1)(b-1)|$ such fixed
    points.
\end{enumerate}

\textit{Part (b).} Note that $T^n(x,y)=(a^nx,b^ny)$. The point $(x_0,y_0)$ of
minimal period $2$ must satisfy $a^2x_0=x_0$ and $b^2y_0=y_0$, implying
$(a^2-1)x_0=(a+1)(a-1)x_0 \in \mathbb{Z}$ and $(b^2-1)y_0=(b-1)(b+1)y_0 \in
\mathbb{Z}$. However, these points cannot be fixed: in other words, such points
are of form $(\frac{n}{a^2-1},\frac{m}{b^2-1})$ where at least one of the
conditions fail: $(a+1)\mid n$ or $(b+1)\mid m$.

We also need to handle special cases when $a = \pm 1$ and $b = \pm 1$. Notice
that when $a=1$, no point can have period $2$ in $x$-coordinate, while for
$a=-1$, every point has period $2$ in $x$-coordinate. Same holds for $b$.

\textit{Part (c).} The system is never minimal. Indeed, if either $a$ or $b$ is
not $1$, then from \textit{part (a)} we know that the system has fixed points.
If $a=b=1$, then $T$ is the identity map which clearly implies the system is not
minimal either.

\textit{Part (d).} We need to understand when the orbit's closure
$\overline{\mathcal{O}(x_0,y_0)}$ of some $(x_0,y_0) \in
\mathbb{R}^2/\mathbb{Z}^2$ is dense. Clearly, if $|a|=|b|=1$, all orbits are
finite and thus not dense, so the system is not transitive. However, when both
$|a|, |b| \geq 2$, then notice that any open set $U$ after applying $T$
stretches and eventually intersects any open set $V$ (thus making $T^n(U) \cap V
\neq \emptyset$). If either of $|a|$ or $|b|$ is $1$, then one of the
coordinates is periodic, and the orbit stays inside closed circles. Thus, the
closure is not dense.

\textbf{Conclusion}: $(X, T)$ is transitive iff $|a| \geq 2$ and $|b| \geq 2$.

\textit{Part (e).} Equivalent to \textit{part (d).}

\textit{Part (f).} The system is weakly mixing whenever $(X \times X, T \times
T)$ is topologically ergodic, or simply transitive in our case. So in other
words, we are given the system $\mathbb{R}^4/\mathbb{Z}^4$ with the map
$T:(x,y,z,w) \mapsto (\alpha x, \beta y, \alpha z, \beta w)$ and we need to
decide whether this system has a transitive point. However, this is clearly the
case as long as $|a|, |b| \geq 2$ (with the similar reasoning to \textit{part
(d)}). So the conclusion is equivalent to \textit{part (d)}.

\textit{Part (g).}  The system is similarly strongly mixing whenever $|a|, |b|
\geq 2$.

\textit{Part (h).} Recall that for $(X,T)$ to be equicontinuous, one has to have
\begin{equation*}
    \forall_{\varepsilon > 0} \; \exists_{\delta > 0} \; d((x_0,y_0),(x_1,y_1))<\delta \Rightarrow d((a^nx_0,b^ny_0), (a^nx_1,b^ny_1))<\varepsilon \; \forall_{n \in \mathbb{Z}_{\geq 0}}
\end{equation*}

However, notice that if either $|a| \geq 2$ or $|b| \geq 2$, then the distance
between two points (in the corresponding coordinate) exponentially increases
(modulo 1). For instance, consider two very close points $(x_0,y_0)$ and
$(x_0+\delta,y_0)$ for some $\delta \ll 1$. After $n$ iterations, the distance
in the first distance becomes $|a|^n\delta$, which grows exponentially (albeit
wrapped around the circle).  

When both $|a|=|b|=1$, then the map $T$ is distance-preserving, thus the 
system is equicontinuous.

\textbf{Conclusion.} $(X,T)$ is equicontinuous iff $|a|=|b|=1$.

\newpage

\begin{problem}[Exercise 2]
    Let $(X,T)$ be a topological dynamical system, and assume that $X$
    is compact and metrizable, with a metric $d$.
    \begin{enumerate}
        \item[(a)] Let $d'$ be another metric on $X$. Show that $(X,T)$ is 
        equicontinuous with respect to $d$ iff $(X,T)$ is equicontinuous with 
        respect to $d'$.
        \item[(b)] Show that $(X,T)$ is equicontinuous iff the following condition 
        is satisfied:
        \begin{equation*}
            \forall_{\varepsilon>0} \; \exists_{K \subseteq \mathbb{N}, K \; \text{finite}} \; \forall_{n \in \mathbb{N}} \; \exists_{k \in K} \; d(T^n(x), T^k(x)) < \varepsilon \; \forall_{x \in X}.
        \end{equation*}
    \end{enumerate}
\end{problem}

\textbf{Solution.}

\textcolor{blue!70!black}{\textbf{Part \textit{(a)}.}} Recall that on a compact 
space, two compatible metrics are uniformly equivalent, i.e.,
\begin{equation*}
    \textcolor{gray}{(1)}: \forall_{\varepsilon>0} \; \exists_{\delta>0} \; \{d'(x,y) < \delta\} \Rightarrow \{d(x,y) < \varepsilon\}.
\end{equation*}

Fix $\varepsilon>0$ and choose $\delta := \varepsilon'$ that satisfies the above inequality. By 
equicontinuity in metric $d$, there exists $\delta>0$ such that
\begin{equation*}
    \textcolor{gray}{(2)}: \{d(x,y) < \delta\} \Rightarrow \{d(T^n(x), T^n(y))<\varepsilon'\} \; \forall_{n \in \mathbb{N}}. 
\end{equation*}

Again by uniform equivalence, choose $\delta'>0$ such that
\begin{equation*}
    \textcolor{gray}{(3)}: \{d'(x,y) < \delta'\} \Rightarrow \{d(x,y) < \delta\}.
\end{equation*}

Then, for any $x,y \in X$, we have
\begin{align*}
    \{d'(x,y)<\delta'\} 
    &\xRightarrow{\textcolor{gray}{(3)}} \{d(x,y)<\delta\} \\
    &\xRightarrow{\textcolor{gray}{(2)}} \{d(T^n(x), T^n(y)) < \varepsilon'\} \; \forall_{n \in \mathbb{N}} \\
    &\xRightarrow{\textcolor{gray}{(1)}} \{d'(T^n(x), T^n(y)) < \varepsilon\} \; \forall_{n \in \mathbb{N}}.
\end{align*}

\textcolor{blue!70!black}{\textbf{Part \textit{(b)}.}} We prove two directions.

\underline{\textit{Equicontinuous $\Rightarrow$ finite covering in $d_{\infty}$.}} Fix some
$\varepsilon > 0$. By equicontinuity, for each $x \in X$ there exists $\delta>0$
such that
\begin{equation*}
    \textcolor{gray}{(1)}: \{d(x,y) < \delta\} \Rightarrow \{d(T^n(x), T^n(y)) < \varepsilon/3\} \; \forall_{n \in \mathbb{N}}.
\end{equation*}

Since $X$ is compact, we can cover $X$ by finitely many balls of radius
$\delta$: concretely, let $X = \bigcup_{i=1}^N B(x_i,\delta)$. However, these
centers on their own do not help much. We need to understand the dynamics of the
sequence $\{(T^n(x_1),\dots,T^n(x_i))\}_{n \in \mathbb{N}}$ which lies in the
space $X^N$ (which further denote by $\{\mathbf{x}_n\}_{n \in \mathbb{N}}$).

That said, we consider the space $(X^N,d_{\infty})$ which we equip with the
metric $d_{\infty}(\mathbf{x},\mathbf{y})=\max_{1 \leq i \leq N}d(x_i,y_i)$.
Clearly, $X^N$ is compact (say, by Tychonoff's Theorem), so the sequence
$\{\mathbf{x}_n\}_{n \in \mathbb{N}}$ is totally bounded. In particular this 
means that we can find a finite set of indices $K \subseteq \mathbb{N}$
such that $\{\mathbf{x}_k\}_{k \in K}$ is the $(\varepsilon/3)$-covering 
of our sequence $\{\mathbf{x}_n\}_{n \in \mathbb{N}}$ in $d_{\infty}$ metric. Specifically:
\begin{equation*}
    \forall_{n \in \mathbb{N}} \; \exists_{k \in K} \; \{d_{\infty}(\mathbf{x}_n, \mathbf{x}_k) < \varepsilon/3\}.
\end{equation*}

So by expanding everything we have so far:
\begin{equation*}
    \textcolor{gray}{(2)}: \exists_{K \subseteq \mathbb{N}, K \text{ is finite}} \; \forall_{n \in \mathbb{N}} \; \exists_{k \in K} \; \{d(T^n(x_i), T^k(x_i)) < \varepsilon/3\} \; \text{for each $i \in \{1,\dots,N\}$}.
\end{equation*}

Now, we have all the tools to finish off the statement. Fix $\varepsilon>0$.
Take $K$ as in \textcolor{gray}{\textit{(2)}}. For every $x \in X$ there exists
some center of the ball $x_i$ such that $d(x,x_i)<\delta$. For every $n \in
\mathbb{N}$ choose $k$ as in \textcolor{gray}{\textit{(2)}}. Then, apply the
triangle inequality:
\begin{equation*}
    d(T^n(x), T^k(x)) \leq 
    \underbrace{d(T^n(x),T^n(x_i))}_{< \varepsilon/3 \, \text{by} \, \textcolor{gray}{\textit{(1)}}} 
    + \underbrace{d(T^n(x_i), T^k(x_i))}_{< \varepsilon/3 \, \text{by} \, \textcolor{gray}{\textit{(2)}}} 
    + \underbrace{d(T^k(x_i), T^k(x))}_{< \varepsilon/3 \, \text{by} \, \textcolor{gray}{\textit{(1)}}}
    < \varepsilon.
\end{equation*}

\underline{\textit{Finite covering in $d_{\infty}$ $\Rightarrow$ Equicontinuous.}} Assume 
\begin{equation*}
    \forall_{\varepsilon>0} \; \exists_{K \subseteq \mathbb{N}, K \; \text{finite}} \; \forall_{n \in \mathbb{N}} \; \exists_{k \in K} \; \forall_{x \in X} \; d(T^n(x), T^k(x)) < \varepsilon.
\end{equation*}

Fix $\varepsilon>0$ and pick any two points $x,y \in X$.

From the fact that $T^k$ is uniformly continuous (in particular for all $k \in
K$ since $X$ is compact), we find $\delta>0$ such that:
\begin{equation*}
    \{d(x,y)<\delta\} \Rightarrow \forall_{k \in K} \; \{\textcolor{red!70!black}{d(T^k(x), T^k(y)) < \varepsilon/3}\}.
\end{equation*}

Then, for any $n$, pick $k \in K$ with $\textcolor{blue!70!black}{d(T^n(x), T^k(x))<\varepsilon/3}$ and 
$\textcolor{green!70!black}{d(T^n(y), T^k(y))<\varepsilon/3}$. Finally, apply the triangle inequality:
\begin{equation*}
    d(T^n(x), T^n(y)) \leq \textcolor{blue!70!black}{d(T^n(x), T^k(x))} + \textcolor{green!70!black}{d(T^k(x), T^k(y))} + \textcolor{red!70!black}{d(T^k(y), T^n(y))} < \varepsilon,
\end{equation*}
hence $(X,T)$ is equicontinuous.

\newpage

\begin{problem}[Exercise 3]
    Consider the following maps on $\mathbb{R}$:
    \begin{gather*}
        T_1(x) = x, \; T_2(x) = -x, \; T_3(x) = x+3, \\ 
        T_4(x) = 4x, \; T_5(x) = 4x(1-x), \; T_6(x) = 6x(1-x).
    \end{gather*}
    Show that no two of the resulting topological dynamical systems $(X,T_i)$ ($1 \leq i \leq 6$) are isomorphic.
\end{problem}

\textbf{Solution.}

\textit{Idea.} To prove that these topological dynamical systems are distinct,
we need to find distinguishing invariant of each system. In this problem 
it suffices to consider periodic points and invariant sets. Formally,
the following proposition states that these two are indeed invariants 
of the dynamical system under isomorphism.

\begin{lemma}[Adapted Proposition 1.5.2, \cite{vries2014}]
    Let $\varphi: (X,T) \xrightarrow{\sim} (Y,S)$ be the isomorphism of dynamical 
    systems. Then,
    \begin{enumerate}
        \item[(a)] point $x$ is periodic in $(X,T)$ iff $\varphi(x)$ is periodic
        in $(Y,S)$ (moreover, primitive periods of both $x$ and $\varphi(x)$ are the same);
        \item[(b)] set $A$ in $(X,T)$ is $T$-invariant iff $\varphi(A)$ is
        $S$-invariant in $(Y,S)$.
    \end{enumerate}
\end{lemma}

\textit{Proof.} By considering periodic points, we can differentiate systems
$T_1,T_2,T_3,T_4$, and (\textit{$T_5$ with $T_6$}) by applying Lemma,
\textit{part (a)}. To differentiate $T_5$ and $T_6$, we are going to look at the
invariant sets and apply Lemma, \textit{part (b)}, although this part 
turned out to be the most challenging in this problem.

\underline{\textit{System $T_1$.}} This is just an identity map, so it is
clearly distinct from other systems: for instance, since every point is fixed.

\underline{\textit{System $T_2$.}} Note that the only fixed points on this map 
is $x=0$ and all other points have the primitive period of $2$. No other 
systems have this characterization.

\underline{\textit{System $T_3$.}} This system, in contrast, has no periodic 
points (all other systems, in turn, have at least a single periodic point).

\underline{\textit{Systems $T_4$.}} This systems has a fixed point $x=0$ and no
other periodic points.

\underline{\textit{Systems $T_5$ and $T_6$.}} This one is the hardest. Let us
consider the general case $f_{\mu}(x) = \mu x(1-x)$ (the so-called
\textit{logistic map}). By solving equation $f_{\mu}(x)=x$ we get that it always
(besides the case $\mu=0$) has two solutions: $x_0=0$ and $x_1=1-1/\mu$.
Although this already differentiates $T_5$ and $T_6$ from the rest of the
systems, we need another measure to differentiate $T_5$ and $T_6$.

We propose to compute the maximum value of $f_{\mu}$ on the set $[0,1] \subseteq
X$. Notice that $f_{\mu}'(x) = \mu(1-2x)$, so the extremum point corresponds to
the value at $x_{\text{ext}} = \frac{1}{2}$ (which lies within $[0,1]$).
Moreover, $f_{\mu}(x_{\text{ext}}) = \frac{1}{4}\mu$. Thus, we've shown the following:

\textbf{Claim \#1.} The set $[0,1]$ is $f_{\mu}$-invariant iff $0 \leq \mu \leq 4$.

Therefore, the dynamical system $T_5$ has an invariant set with non-empty
interior. Consequently, by Lemma, \textit{part (b)}, if there were the
isomorphism between $(X,T_5)$ and $(X,T_6)$, it would correspond to the
invariant set with non-empty interior in $(X,T_6)$. However, our second claim 
is as follows.

\textbf{Claim \#2.} The dynamical system $(X,f_{\mu})$ has no invariant sets with non-empty
interior for $\mu > 2+\sqrt{5}$.

\textit{Proof of claim.} First we observe that the orbit of any point outside
$[0,1]$ ``explode'' to $-\infty$ (for instance, take any $x_0>1$, then $x_1 :=
f_{\mu}(x_0)$ becomes negative, and the next iteration will scale $x_1$ by the
factor $\mu(1-x_1)$, which is clearly more than $1$, and thus will repel the
point further to $-\infty$). 

Therefore, if we had an invariant set $A \subseteq X$, then the orbit
$\mathcal{O}(x)$ of any $x \in A$ must remain in $[0,1]$. Let us check whether
this holds.

Let us check the ``escape region'' $E \subseteq X$, from which we get values
greater than $1$: i.e., $f_{\mu}(x) > 1$. This yields equation $\mu x^2 - \mu x
+ 1 < 0$ which gives the interval $E = (\ell,r)$ for 
\begin{equation*}
    \ell = \frac{1}{2} - \frac{1}{2}\sqrt{1-\frac{4}{\mu}}, \quad r = \frac{1}{2} + \frac{1}{2}\sqrt{1-\frac{4}{\mu}}.
\end{equation*} 

\textcolor{gray}{\textit{(Note that this is a valid non-empty interval as $\mu>4$)}.}

Therefore, the ``safe regions'' of $f_{\mu}$ are two zones: $L=[0,\ell]$ and
$R=[r,1]$: if the point happens to be in $(\ell,r)$, it will be sent to
infinity.

Now take any open interval $(a,b) \subseteq A$. Without loss of generality,
assume $(a,b) \subseteq L$ or $(a,b) \subseteq R$ (otherwise, if $E \cap (a,b)
\neq \emptyset$, then clearly $A$ is not invariant). We claim that $f_{\mu}^n((a,b))$
eventually intersects $E$. Indeed, it suffices to verify that this interval gets
expanded with each iteration of $f_{\mu}$. To check this, we check the minimum value
of $|f_{\mu}'|$ in the ``safe'' zone $L \cup R$\footnote{Recall that from Mean
Value Theorem, $f(x_2)-f(x_1) = f'(c)(x_2-x_1)$ for some $c \in [x_1,x_2]$.}. 

Note that since $f_{\mu}'(x)$ is a linear function, its magnitude strictly
decreases at $x$ gets closer to the critical point $x_{\text{ext}}=\frac{1}{2}$.
Thus, the minimum value of $|f_{\mu}'(x)|$ on $L \cup R$ occurs at points $\ell$
and $r$. One easily checks $|f_{\mu}'(\ell)| = |f_{\mu}'(r)| =
\sqrt{\mu^2-4\mu}$. If this value exceeds $1$, then the claim is proven. 
This occurs when $\mu^2-4\mu-1>0$ which reduces to $\mu>2+\sqrt{5}$.

\textbf{Corollary.} $(X,T_6)$ has no invariant sets with non-empty interior.

\textbf{Answer.} See \Cref{tab:system_comparison} below.

\begin{table}[htpb]
    \centering
    \renewcommand{\arraystretch}{1.3}
    \begin{tabular}{cll}
        \hline
        \textbf{System} & \textbf{Periodic Points} & \textbf{Invariant Sets $A$ ($\text{int}(A) \neq \emptyset$)} \\
        \hline
        $T_1$ & Every point is a fixed point & N/A \\
        \hline
        $T_2$ & $x=0$ is fixed; others have min. period 2 & N/A \\
        \hline
        $T_3$ & None & N/A \\
        \hline
        $T_4$  & $x=0$ is fixed; no others & N/A \\
        \hline
        $T_5$  & Two fixed points ($x=0, x=3/4$) & Yes ($[0,1]$ is invariant) \\
        \hline
        $T_6$  & Two fixed points ($x=0, x=5/6$) & None \\
        \hline
    \end{tabular}
    \caption{Comparison of distinguishing invariants for dynamical systems $(X, T_i)$.}
    \label{tab:system_comparison}
\end{table}

\newpage

\begin{problem}[Exercise 4]
    We call a topological dynamical system $(X,T)$ \textit{totally minimal}
    (resp. \textit{totally transitive}) if $(X,T^d)$ is minimal (resp. transitive)
    for all $d \geq 1$. 
    \begin{enumerate}
        \item[(a)] Show that for $\alpha \in \mathbb{R} \setminus \mathbb{Q}$, the 
        circle rotation by $\alpha$ is totally minimal.
        \item[(b)] Show that the doubling map on the circle is totally transitive.
        \item[(c)] Show that not each minimal (resp. transitive) topological
        dynamical system is totally minimal (resp. totally transitive). 
    \end{enumerate}
\end{problem}

\textbf{Solution.}

\textcolor{blue!70!black}{\textbf{Part \textit{(a)}.}} It is well-known that the
map $R_{\alpha}: x \mapsto x+\alpha$ (over $\mathbb{R}/\mathbb{Z}$) is minimal
iff $\alpha \in \mathbb{R}\setminus \mathbb{Q}$. Note that
$R_{\alpha}^d=R_{d\alpha}$ and $d\alpha$ is clearly irrational as well, thus
$R_{\alpha}^d$ is also minimal. Hence, $(X,R_{\alpha})$ is totally minimal. 

\textcolor{blue!70!black}{\textbf{Part \textit{(b)}.}} We already know well that
the doubling map is transitive. In fact, any map $M_{\gamma}: x \mapsto \gamma
x$ for integer $\gamma \geq 2$ is transitive. The argument is similar to the
case when $\gamma=2$: we need to show that for any opens $U,V \subseteq
\mathbb{R}/\mathbb{Z}$ there exists $n \in \mathbb{N}$ such that
$T_{\gamma}^n(U) \cap V \neq \emptyset$. If $U$ is an open interval of length 
$\delta>0$, then after $n$ iterations the ``length'' of $T_{\gamma}^n(U)$ is 
$\gamma^n\delta$. For large enough $n$ this exceeds $1$ and so $T_{\gamma}^n(U)$
effectively wraps the original $U$ around the circle at least once. Thus,
as we cover the whole circle, we definitely intersect $V$ at some point.

That said, the $n$th iterate of the doubling map satisfies $M_2^n=M_{2^n}$ which
is transitive by the argument above for every $n$. Thus, $(X,M_2)$ is totally
transitive.

\textcolor{blue!70!black}{\textbf{Part \textit{(c)}.}}

\textit{Minimal but not totally minimal.} Let $X=(\mathbb{R}/\mathbb{Z}) \times
\mathbb{Z}_2$ (where $\mathbb{Z}_2$ is equipped with discrete topology) and $T:
(x,b) \mapsto (x+\alpha,b+1)$ (with $\alpha \in \mathbb{R}/\mathbb{Z}$).
Clearly, $T$ is minimal, but is not totally minimal since
$T^2(x,b)=(x+2\alpha,b)$ is not minimal as it fixes the second coordinate (thus,
say, $(\mathbb{R}/\mathbb{Z}) \times \{0\}$ is a smaller $T^2$-invariant set).

\textit{Transitive but not totally transitive.} Similarly consider
$X=(\mathbb{R}/\mathbb{Z}) \times \mathbb{Z}_2$ and $T: (x,b) \mapsto (2x,b+1)$.
Then, $T$ is transitive while $T^2(x,b)=(4x,b)$ is not since the only possible
orbits are $(\mathbb{R}/\mathbb{Z}) \times \{0\}$ or $(\mathbb{R}/\mathbb{Z})
\times \{1\}$.

\newpage

\begin{problem}[Exercise 5]
    Let $\phi: X \to Y$ be a morphism between topological dynamical systems 
    $(X,T)$ and $(Y,S)$. Assume that $X$ is locally compact, and let $A \subseteq X$
    be non-empty, compact and invariant.
    \begin{enumerate}
        \item[(a)] Show that $\phi(\mathcal{B}_T(A)) \subseteq \mathcal{B}_S(\phi(A))$,
        where $\mathcal{B}_T(A)$ denotes the basin of attraction of the set $A$ under 
        the map $T$.
        \item[(b)] Show that the inclusion in the previous point can be strict.
    \end{enumerate}
\end{problem}

\textbf{Solution.}

\textcolor{blue!70!black}{\textbf{Part \textit{(a)}.}} Recall that the basin of
attraction is by definition 
\begin{equation*}
    \mathcal{B}_T(A) = \{x \in X: \emptyset \neq \omega_T(x) \subseteq A\}.  
\end{equation*}

We prove this fact by the series of three claims.
\begin{enumerate}[label=\textbf{Claim \#\arabic*}, align=left]
    \item Assume $X$ is locally compact. If $x \in X$ and $\omega(x)$ is compact and not empty, then $\overline{\mathcal{O}}(x)$ is compact.
    \item If $\overline{\mathcal{O}}(x)$ is compact then $\phi(\omega_T(x)) = \omega_S(\phi(x))$.
    \item Assume $\phi(A)$ is a closed subset of $Y$ and all
    $\overline{\mathcal{O}}(x)$ for $x \in \mathcal{B}_T(A)$ are compact. Then, $\phi(\mathcal{B}_T(A)) \subseteq \mathcal{B}_S(\phi(A))$.
\end{enumerate}

\textbf{Proof of claims.}

\textit{Claim \#1.} Since $X$ is locally compact, $\omega(x)$ has an open
neighbourhood $U$ with compact closure $\overline{U}$. Note that $T^n(x) \in U
\subseteq \overline{U}$ holds for infinitely many values of $n$. Thus, it
suffices to show that $T^n(x) \in \overline{U}$ for almost all values of $n$.
Assume, however, that it is not the case. Then, there is a sequence
$\{T^{n_i}(x)\}_{i \in \mathbb{N}} \subseteq \overline{U}$ such that
$T^{n_i+1}(x) \not\in \overline{U}$ for all $i$. The first sequence has an
accumulation point $z$ in $\omega(x)$, but the second sequence has an
accumulation point $T(z)$ outside $U$, contradicting invariance of $\omega(x)$.

\textit{Claim \#2.} Note that in general we have $\phi(\omega_T(x)) \subseteq \omega_S(\phi(x))$:
\begin{align*}
    \phi(\omega_T(x)) &\triangleq \phi\left(\bigcap_{n \geq 0}\overline{\mathcal{O}}_T(T^n(x))\right) & \textcolor{gray}{\textit{// definition of limit set}} \\
    & \subseteq \bigcap_{n \geq 0}\phi(\overline{\mathcal{O}}_T(T^n(x))) & \textcolor{gray}{\textit{// use the fact that $\phi$ is continuous}} \\
    & \subseteq \bigcap_{n \geq 0}\overline{\phi(\mathcal{O}_T(T^n(x)))} & \textcolor{gray}{\textit{// again, since $\phi$ is continuous}} \\
    & = \bigcap_{n \geq 0}\overline{\mathcal{O}}_S(\phi(T^n(x))) & \textcolor{gray}{\textit{// just a direct observation}} \\
    & = \bigcap_{n \geq 0}\overline{\mathcal{O}}_S(S^n(\phi(x))) & \textcolor{gray}{\textit{// definition of morphism}} \\
    & = \omega_S(\phi(x)) & \textcolor{gray}{\textit{// definition of limit set}} \\
\end{align*}

Note that we remove inclusions at second and third step using the fact that
$\overline{\mathcal{O}}(x)$ is compact.

\textit{Claim \#3.} By \textit{Claim \#2}, $\omega_S(\phi(x)) =
\phi(\omega_T(x))$ is a non-empty subset of $\phi(A)$ for all $x \in
\mathcal{B}_T(A)$, hence $\phi(\mathcal{B}_T(A)) \subseteq
\mathcal{B}_S(\phi(A))$.

Now simply apply three claims. Since $A$ is compact, then for every point $x \in
\mathcal{B}_T(A)$, the set $\omega(x)$ is non-empty and compact. By
\textit{Claim \#1}, $\overline{\mathcal{O}}(x)$ is compact, so by \textit{Claim
\#3} we have proven the theorem.

\textcolor{blue!70!black}{\textbf{Part \textit{(b)}.}} Take the following
example:

\begin{itemize}
    \item \underline{System $(X,T)$.} Let $X = \{0\} \cup [2,3]$ and let
    \begin{equation*}
        T(x) = \begin{cases}
            0, & x = 0, \\
            1+\frac{1}{2}x, & x \in [2,3].
        \end{cases}
    \end{equation*}
    Set $A=\{0\}$ (so that $\mathcal{B}_T(A)=\{0\}$).
    \item \underline{System $(Y,S)$.} Let $Y=[0,1]$ and $S(y) = y/2$.
\end{itemize}

Let the morphism $\phi: (X,T) \to (Y,S)$ be 
\begin{equation*}
    \phi(x) = \begin{cases}
        0, & x=0, \\
        x-2, & x \in [2,3].
    \end{cases}
\end{equation*}

This is a valid morphism: it is continuous and satisfies $\phi(T(x)) =
S(\phi(x))$.

Finally, note:
\begin{itemize}
    \item $\phi(\mathcal{B}_T(A)) = \phi(\{0\}) = \{0\}$.
    \item $\mathcal{B}_S(\phi(A)) = \mathcal{B}_S(\{0\}) = [0,1]$ as any point
    of $Y$ converges to $0$.
\end{itemize}

Thus, we see that $\{0\}=\phi(\mathcal{B}_T(A)) \subsetneq \mathcal{B}_S(\phi(A)) = [0,1]$.

%\nocite{*}
%\renewcommand{\refname}{}
\bibliographystyle{alphaurl}
\bibliography{refs}

\end{document}